\article{$\widetilde{\operatorname{SL}}_2$ 尖点形式与正交群自守形式之间的一种关系}{On a Relation between $\widetilde{\operatorname{SL}}_2$ Cusp Forms and Automorphic Forms on Orthogonal Groups}{S. Rallis}{1}{297--314}

本讲从 A. Weil 在 [22] 中的思想出发，说明如何在 $\widetilde{\operatorname{SL}}_2$ 尖点形式（cusp form）与正交群上的自守形式（automorphic form）之间构造一个对应或提升（lifting）。

\origsection{1}{Siegel 公式与模形式的提升（Siegel formula and the lifting of modular forms）}
设 $\R^k$ 为 $k$ 维欧几里得空间，标准基为 $e_1,\ldots,e_k$。以 $[\ ,\ ]$ 表示 $\R^k$ 上的双线性型
\[
 [X,Y]=\sum_{i=1}^k x_i y_i,\qquad X=\sum_{i=1}^k x_i e_i,\quad Y=\sum_{i=1}^k y_i e_i.
\]
在 $\R^k$ 上定义双线性型 $Q$：
\[
 Q(X,Y)=\sum_{i=1}^a x_i y_i-\sum_{i=a+1}^k x_i y_i,\qquad a<k,
\]
其对应的二次型（quadratic form）满足：对每个 $X\in\R^k$，
\[
 Q(X,X)=\|X_+\|^2-\|X_-\|^2,\quad
 \|X_+\|^2=\sum_{i=1}^a x_i^2,\quad\|X_-\|^2=\sum_{i=a+1}^k x_i^2.
\]
因此，$Q$ 的符号（signature）为 $(a,b)$，其中 $a+b=k$ 且 $b\ge1$。

若格（lattice）$L\subseteq\R^k$ 满足 $Q(L,L)\subseteq\Z$，则称之为 $Q$ 整的（$Q$-integral）；若 $Q(L,L)\subseteq2\Z$，则称之为 $Q$ 偶的（$Q$-even）。$L$ 的 $Q$ 对偶格（$Q$-dual lattice）为
\[
 L_*(Q)=\{\xi\in\R^k:Q(\xi,L)\subseteq\Z\}.
\]
那么 $L_*(Q)/L$ 是有限 Abel 群。以 $n_L$ 表示这个群的指数（exponent），即使得对所有 $\xi\in L_*(Q)$ 均有 $n_L\xi\in L$ 的最小正整数。若 $\xi_1,\ldots,\xi_k$ 是 $L$ 的任意一组 $\Z$ 基，令
\[
 D_{Q(L)}=\det\{Q(\xi_i,\xi_j)\}.
\]
这个数称为 $Q$ 相对于 $L$ 的判别式（discriminant），与所选的 $\Z$ 基无关。

定义函数
\begin{equation}
 \psi(X,z,g)=e^{\pi i\{xQ(X,X)+iy[gX,gX]\}}(\operatorname{Im}z)^{b/2},\tag{1-1}
\end{equation}
其中 $z=x+iy\in H$，$H$ 是上半平面（即 $y>0$），而 $g\in O(Q)$，后者为 $Q$ 的正交群（orthogonal group）。取 $\eta\in L_*(Q)$，考虑 theta 级数（theta-series）
\begin{equation}
 \widetilde\Theta_{\psi,\eta}^L(z,g)=\sum_{\xi\in L}\psi(\xi+\eta,z,g).\tag{1-2}
\end{equation}
令 $\eta$ 遍历 $L_*(Q)/L$ 的一组代表元，组成列向量
\begin{equation}
 \Psi_\psi^L(z,g)=\bigl(\widetilde\Theta_{\psi,\eta}^L(z,g)\bigr)_{\eta\in L_*(Q)/L}.\tag{1-3}
\end{equation}
在 $L$ 为 $Q$ 偶格的条件下，Siegel 证明了如下函数方程：
\begin{equation}
 \Psi_\psi^L(\Gamma\cdot z,g\gamma)
 =c(\Gamma,\gamma)(c_\Gamma z+d_\Gamma)^{(a-b)/2}\Psi_\psi^L(z,g),\tag{1-4}
\end{equation}
其中 $\Gamma\in\operatorname{SL}_2(\Z)$，
\[
 \gamma\in O(Q)_L=\{g\in O(Q):g(L)=L\},
\]
$c(\Gamma,\gamma)$ 是一个仅依赖于 $\Gamma,\gamma$、阶数为 $|L_*(Q)/L|$ 的酉矩阵（unitary matrix）。一般而言，映射 $(\Gamma,\gamma)\mapsto c(\Gamma,\gamma)$ 定义群 $\widetilde{\operatorname{SL}}_2(\Z)\times O(Q)_L$ 的一个射影酉表示（projective unitary representation）。

研究 $\Psi_\psi^L(z,g)$ 的主要困难在于，这个函数并非自守的，因为它不是 $\widetilde{\operatorname{SL}}_2\times O(Q)$ 的包络代数（enveloping algebra）的中心的特征函数（eigenfunction）。将 $\widetilde\Theta_{\psi,\eta}^L(z,g)$ 与 $O(Q)/O(Q)_L$ 上的常数函数配对积分，可以解决这一问题。具体来说，Siegel 公式的分析部分是在齐性空间（homogeneous space）
\[
 O(Q)/\{\gamma\in O(Q)_L:\gamma\eta\equiv\eta\pmod L\}
 \simeq O(Q)/O(Q)_{L,\eta}
\]
上选取适当归一化的 $O(Q)$ 不变测度 $d\mu_\eta$，从而得到恒等式
\begin{equation}
 \int_{O(Q)/O(Q)_{L,\eta}}\widetilde\Theta_{\psi,\eta}^L(z,g)\,d\mu_\eta(g)
 =E_\eta(z),\tag{1-5}
\end{equation}
该式在 $k\ge5$ 且 $L$ 为 $Q$ 偶格时成立 [20]，其中
\begin{equation}
\begin{split}
 E_\eta(z)={}&(\operatorname{Im}z)^{b/2}
 \sum_{M\in\operatorname{SL}_2(\Z)_\infty\backslash\operatorname{SL}_2(\Z)}
 c_{\eta,0}^{-}((M,1))\\[-2pt]
 &\hspace{25mm}\cdot(c_Mz+d_M)^{(b-a)/2}|c_Mz+d_M|^{-b}.
\end{split}\tag{1-6}
\end{equation}
\footnote{译注：扫描本在系数 $c_{\eta,0}$ 后显示一个短横，原文未单独解释其含义；这里以 $c_{\eta,0}^{-}$ 保留该记号。}
这里
\[
 \operatorname{SL}_2(\Z)_\infty=\{\Gamma\in\operatorname{SL}_2(\Z):c_\Gamma=0\},
\]
$c_{\eta,0}^{-}((M,1))$ 表示矩阵 $c((M,1))$ 第一列中的一个元素，而 $1$ 为 $O(Q)$ 的单位元素。我们知道 $E_\eta(z)$ 是 $\widetilde{\operatorname{SL}}_2$ 的包络代数的中心的特征函数；具体地说，它满足方程\footnote{\textbf{译者注（本次校订）：}式 (1-7) 与随后所列的微分算子不相容。例如 (1-6) 中恒等陪集给出的项为 $y^{b/2}$，代入所列算子得到 $\tfrac b2(\tfrac b2+\tfrac k2-1)y^{b/2}$，并不为零。按 (1-4) 的权 $w=(a-b)/2$，相应的正确方程为 $[y^2(\partial_x^2+\partial_y^2)+wy\partial_y-iwy\partial_x-b(a-2)/4]E_\eta=0$；该算子逐项作用于 (1-6) 即可验证。}
\begin{equation}
 \Delta_{k/2}(E_\eta)=0,\tag{1-7}
\end{equation}
其中
\[
 \Delta_{k/2}=y^2\left(\frac{\partial^2}{\partial x^2}+\frac{\partial^2}{\partial y^2}\right)
 +\tfrac12ky\frac{\partial}{\partial y}
 +\tfrac i2(b-a)y\frac{\partial}{\partial x}.
\]
用表示论的语言说，Siegel 公式表达了如下事实：$O(Q)$ 的平凡表示，在下文将精确说明的意义下，对应于由非解析 Eisenstein 级数 $(E_\eta)_\eta$ 决定的 $\widetilde{\operatorname{SL}}_2$ 的表示。这一构造还可推广到 $O(Q)$ 的其他自守表示，具体分为两步。首先改变核函数（kernel function）。令 $P$ 为 $\R^a$（即 $\{e_1,\ldots,e_a\}$ 的实线性张成空间）上次数为 $t$ 的齐次多项式（homogeneous polynomial），满足 $\partial(Q)(P)=0$，其中 $\partial(Q)$ 为微分算子 $Q(\partial/\partial x_1,\ldots,\partial/\partial x_n)$。构造
\[
 \psi_*(X,z,g)=(\operatorname{Im}z)^{t/2}
 P\bigl(g\cdot(\operatorname{Im}z)^{1/2}X\bigr)\psi(X,z,g),
\]
再如上定义 $\Psi_{\psi_*}^L(z,g)$。于是 $\Psi_{\psi_*}^L$ 满足类似 (1-4) 的函数方程。第二步，是将 $\widetilde\Theta_{\psi_*,\eta}^L$ 与 $O(Q)/O(Q)_{L,\eta}$ 上的尖点形式 $\beta$ 相乘后积分。这是可行的，因为 $\widetilde\Theta_{\psi_*,\eta}^L$ 是 $H\times O(Q)/O(Q)_{L,\eta}$ 上的“缓慢增长”函数。

以
\[
 \langle\widetilde\Theta_{\psi_*,\eta}^L(z,g)\mid\beta(g)\rangle_{O(Q)/O(Q)_{L,\eta}}
\]
表示 $\widetilde\Theta_{\psi_*,\eta}^L$ 与 $\beta$ 关于 $d\mu_\eta$ 的 Hilbert 空间内积。若 $\omega_{O(Q)}*\beta=\lambda\beta$，其中 $\omega_{O(Q)}$ 为 $O(Q)$ 的 Casimir 元（Casimir element），则可推出
\[
 z\longmapsto\langle\widetilde\Theta_{\psi_*,\eta}^L(z,g)\mid\beta(g)\rangle
\]
是一个类似 $\Delta_{k/2}$ 的算子的特征函数。至于所得函数是否为尖点形式，则需要另行证明。文献 [19] 在 $O(2,1)$ 的情形证明：若 $\beta$ 是 $\operatorname{PGL}_2(\R)$（至少局部同构于 $O(2,1)$）上的“全纯”尖点形式，则对 $\operatorname{SL}_2(\Z)$ 的一个适当的 Hecke 同余子群（Hecke congruence subgroup）（取决于格 $L$），上述函数是在 [18] 意义下的半整权全纯尖点形式。文献 [9] 则考察了 $O(1,1)$、各向异性格 $L$ 以及 $O(Q)/O(Q)_{L,\eta}$ 上的酉特征 $\beta$；在这一情形，上述函数对应于一个由 $\operatorname{SL}_2(\R)$ 的主系列表示（principal series representation）决定的非解析尖点形式（亦见 [8]）。

另一方面，也可尝试用类似的构造反转这一对应。令 $h(z)$ 是 $\widetilde{\operatorname{SL}}_2$ 上的尖点形式，仍采用 [17] 中的定义，并取 $\operatorname{SL}_2(\Z)$ 的一个适当 Hecke 同余子群 $\Gamma$。构造 $h$ 与 $\widetilde\Theta_{\psi_*,\eta}^L$ 的 Petersson 内积（Petersson inner product）
\begin{equation}
 \langle\widetilde\Theta_{\psi_*,\eta}^L(\ ,g)\mid h(\ )\rangle
 =\int_{\Gamma\backslash H}\widetilde\Theta_{\psi_*,\eta}^L(z,g)
 \overline{h(z)}(\operatorname{Im}z)^{\nu-2}\,dx\,dy,\tag{1-8}
\end{equation}
其中 $\nu$ 是 $h$ 的权。于是容易推出
\[
 g\longmapsto\langle\widetilde\Theta_{\psi_*,\eta}^L(\ ,g)\mid h(\ )\rangle
\]
是 $O(Q)$ 的包络代数的中心的特征函数。所得函数是否为尖点形式，同样需要另行讨论。Niwa [10] 对 $O(2,1)$ 的情形说明，这种对应如何解释 Shimura [18] 从半整权尖点形式构造整权尖点形式的方法。

因此，研究的核心是核函数 $\widetilde\Theta_{\psi_*,\eta}^L$。我们的思路是，通过详细研究 Weil 表示（Weil representation）的谱分解（spectral decomposition），可以分析 $\widetilde\Theta_{\psi_*,\eta}^L$ 的“尖点分量”，即它在上述各类尖点形式空间上的“投影”。确定这些“尖点分量”后，便能建立所讨论两个群上的模形式之间的对应或提升。

\origsection{2}{Weil 表示的离散谱（Discrete spectrum of the Weil representation）}
定义 Fourier 变换 $\widehat{\ }$ 如下：若 $f\in L^1\cap L^2$，则
\[
 \widehat f(Y)=\int_{\R^k}f(X)e^{2\pi i[X,Y]}\,dX.
\]
令 $\pi_Q$ 为 $\widetilde{\operatorname{SL}}_2\times O(Q)$ 在 $L^2(\R^k)$ 中的 Weil 表示，其中 $\widetilde{\operatorname{SL}}_2$ 为 $\operatorname{SL}_2(\R)$ 的二重覆盖。特别地，若
\[
 \left(\left(\begin{pmatrix}a&b\\0&a^{-1}\end{pmatrix},1\right),g\right)
 \in\widetilde{\operatorname{SL}}_2\times O(Q),\qquad a>0,
\]
则
\[
 \pi_Q\!\left(\left(\begin{pmatrix}a&b\\0&a^{-1}\end{pmatrix},1\right),g\right)\varphi(Z)
 =|a|^{k/2}e^{\pi iabQ(Z,Z)}\varphi(ag^{-1}\cdot Z).
\]
又若
\[
 w_0=\left(\left(\begin{pmatrix}0&-1\\1&0\end{pmatrix},-1\right),1\right),
\]
则 $\pi_Q(w_0)\varphi(Z)=\delta_Q\widehat\varphi(-M_Q(Z))$，其中 $M_Q\in\operatorname{Aut}(\R^k)$ 满足对所有 $X,Y\in\R^k$，$Q(X,Y)=[X,M_Q(Y)]$，而 $\delta_Q$ 为某个八次单位根。上述公式决定了 $\pi_Q$。

令 $\widetilde K$ 为 $\widetilde{\operatorname{SL}}_2$ 的极大紧子群（maximal compact subgroup）：
\[
 \left\{k(\theta,\varepsilon)=
 \left(\begin{pmatrix}\cos\theta&\sin\theta\\-\sin\theta&\cos\theta\end{pmatrix},\varepsilon\right)
 :-\pi\le\theta<\pi,\ \varepsilon=\pm1\right\}.
\]
并令
\[
 A=\left\{a(r)=\left(\begin{pmatrix}r&0\\0&r^{-1}\end{pmatrix},1\right):r>0\right\},
\quad
 N=\left\{n(x)=\left(\begin{pmatrix}1&x\\0&1\end{pmatrix},1\right):x\in\R\right\}.
\]
若 $\mathfrak a,\mathfrak n,\mathfrak t$ 分别为 $A,N,\widetilde K$ 的无穷小生成元（infinitesimal generator），则
\[
 \omega_{\operatorname{SL}_2}=-\mathfrak t^2+\mathfrak a^2+
 (\mathfrak n+\operatorname{Ad}(w_0)\mathfrak n)^2.
\]
令 $K\simeq O(a)\times O(b)$ 为 $O(Q)$ 的极大紧子群。若 $\pi$ 是 $\widetilde{\operatorname{SL}}_2\times O(Q)$ 在 Banach 空间 $\mathcal B$ 上的可微表示，则以 $\mathcal B_{\widetilde K\times K}$ 表示 $\mathcal B$ 中的 $\widetilde K\times K$ 有限向量（finite vector）空间。

令 $F_Q$ 为 $\pi_Q$ 在 $L^2(\R^k)$ 中的 $C^\infty$ 向量空间。由 Sobolev 正则性定理（Sobolev regularity theorem）可知，若 $\varphi\in F_Q$，则对包络代数 $u(\widetilde{\operatorname{SL}}_2\times O(Q))$ 的任意元素 $\gamma$，$\pi_Q(\gamma)\varphi$ 都是 $\Omega_+$ 和 $\Omega_-$ 上的 $C^\infty$ 函数，其中
\[
 \Omega_+=\{X\in\R^k:Q(X)>0\},\qquad
 \Omega_-=\{X\in\R^k:Q(X)<0\}.
\]
用不变量理论的一个简单计算可验证，$\widetilde{\operatorname{SL}}_2$ 与 $O(Q)$ 的包络代数的中心在 $\pi_Q$ 中具有相同的像。即
\[
 \pi_Q(\mathfrak z(\widetilde{\operatorname{SL}}_2))=\pi_Q(\mathfrak z(O(Q))),
\]
其中 $\mathfrak z(\ )$ 表示相应包络代数的中心。特别地，有
\[
 \pi_Q(\omega_{\operatorname{SL}_2})+\pi_Q(\omega_{O(Q)})
 =(\tfrac14k^2-k)I,
\]
这里两个 $\omega$ 分别是 $\widetilde{\operatorname{SL}}_2$ 与 $O(Q)$ 的 Casimir 元。因此，$L^2(\R^k)$ 中代数 $\pi_Q(\mathfrak z(\widetilde{\operatorname{SL}}_2\times O(Q)))$ 的谱分析，归结为单个算子 $\pi_Q(\omega_{\operatorname{SL}_2})$ 的谱分析。

\textit{本文以下始终假设 $k\ge3$ 且 $a\ge b\ge1$。}

令
\[
 F_Q(\lambda)=\{\varphi\in F_Q:\pi_Q(\omega_{\operatorname{SL}_2})\varphi=\lambda\varphi\},
 \quad\operatorname{Spec}(\pi_Q)=\{\lambda\in\C:F_Q(\lambda)\ne\{0\}\}.
\]
因此，若 $\pi$ 是 $\widetilde{\operatorname{SL}}_2\times O(Q)$ 的一个作为离散直和项出现在 $L^2(\R^k)$ 中的不可约酉表示，则 $\pi$ 是某个 $F_Q(\lambda)$ 的子表示。更一般地有如下结论。

\begin{originaltheorem}{定理 2.1（Theorem 2.1）[12]、[14]}
(a) $\operatorname{Spec}(\pi_Q)=\{s^2-2s:s\equiv\tfrac12k\pmod1,\ s>1\}$。

(b) 令
\[
 F_Q^+(\lambda)=\{\varphi\in F_Q(\lambda):\varphi|_{\Omega_-}=0\},\quad
 F_Q^-(\lambda)=\{\varphi\in F_Q(\lambda):\varphi|_{\Omega_+}=0\}.
\]
那么 $F_Q^+(\lambda),F_Q^-(\lambda)$ 是拓扑不可约（topologically irreducible）的 $\widetilde{\operatorname{SL}}_2\times O(Q)$ 模，并有直和
\[
 F_Q(\lambda)=F_Q^+(\lambda)\oplus F_Q^-(\lambda).
\]
此外，$F_Q^+(\lambda)$ 与 $F_Q^-(\lambda)$ 是不等价的不可约表示。\footnote{\textbf{译者注（本次校订）：}本项的不可约性应限于非零分量，不能断言两者总是非零。例如 $b=1$ 时，$F_Q^-(s^2-2s)=\{0\}$；第4节引理4.1之后也明确使用了这一事实。因而这里的分解容许一个分量为零。}

(c) $F_Q^+(s^2-2s)$ 作为 $\widetilde{\operatorname{SL}}_2\times O(Q)$ 表示等价于张量积
\[
 L^2(\mathrm{Whit})_{s^2-2s}\widehat\otimes\mathcal A_s^+,
\]
其中 $L^2(\mathrm{Whit})_{s^2-2s}$ 是从酉特征
\[
 \left(\begin{pmatrix}1&x\\0&1\end{pmatrix},1\right)\longmapsto e^{\pi ix}
\]
酉诱导得到的 $\widetilde{\operatorname{SL}}_2$ 表示的 $C^\infty$ 向量空间中，$\omega_{\operatorname{SL}_2}$ 的特征值为 $s^2-2s$ 的特征空间；$\mathcal A_s^+$ 是 $O(Q)$ 在 $L^2(\Gamma_{+1},d\mu_{+1})$ 中的标准表示的 $C^\infty$ 向量空间中，Laplace--Beltrami 算子 $W_\xi^+$ 的特征值为 $s^2-2s-\tfrac14k^2+k$ 的特征空间。这里
\[
 \Gamma_{+1}=\{X\in\R^k:Q(X,X)=1\},
\]
$d\mu_{+1}$ 是 $\Gamma_{+1}$ 上的 $O(Q)$ 不变测度；$W_\xi^+$ 由 $\Omega_+$ 中的变量分离确定：
\[
 \partial(Q)=\frac{\partial^2}{\partial t^2}
 +\frac{k-1}{t}\frac{\partial}{\partial t}-\frac1{t^2}W_\xi^+,
 \qquad X=t\xi,\ \xi\in\Gamma_{+1}.
\]
\end{originaltheorem}

\paragraph*{注 2.1（Remark 2.1）.}
$F_Q^-(s^2-2s)$ 作为 $\widetilde{\operatorname{SL}}_2\times O(Q)$ 表示，等价于张量积
\[
 L^2(\mathrm{Whit})_{s^2-2s}^{A}\widehat\otimes\mathcal A_s^-,
\]
其中 $L^2(\mathrm{Whit})_{s^2-2s}^{A}$ 是 $L^2(\mathrm{Whit})_{s^2-2s}$ 上的 $\widetilde{\operatorname{SL}}_2$ 表示被 $\widetilde{\operatorname{SL}}_2$ 的唯一外自同构（outer automorphism）$A$ 扭曲后所得的表示，而 $\mathcal A_s^-$ 是 $O(Q)$ 在 $L^2(\Gamma_{-1},d\mu_{-1})$ 中的标准表示的 $C^\infty$ 向量空间中，Laplace--Beltrami 算子 $W_\xi^-$ 的特征值为 $s^2-2s-\tfrac14k^2+k$ 的特征空间。

分析 $\pi_Q$ 的离散谱时，首先要观察到：$F_Q(\lambda)$ 中的一个 $\widetilde K$ 特征函数 $\varphi$（也假设它是 $K$ 有限的）在 $\Omega_\pm$ 上有变量分离性质（separation of variables），即
\[
 \varphi(X)=\rho_1^\pm(t)\rho_2^\pm(\xi),
\]
其中径向函数 $\rho_1^\pm$ 与角向函数 $\rho_2^\pm$ 分别满足相应的微分方程，且 $X=t\xi$，$\xi\in\Gamma_{\pm1}$，$t\in\R_+$。其次，映射
\[
 \varphi\longmapsto\bigl[G\longmapsto\pi_Q(G^{-1})(\varphi)(\xi),
 \quad G\in\widetilde{\operatorname{SL}}_2\bigr]
\]
定义了从 $F_Q(\lambda)$ 中的 $\widetilde K$ 有限函数到从酉特征
\[
 n(x)\longmapsto e^{\pi ix\operatorname{sgn}Q(\xi,\xi)}
\]
酉诱导得到的 $\widetilde{\operatorname{SL}}_2$ 表示的“平方可积”函数空间的一个 $\widetilde{\operatorname{SL}}_2$ 无穷小交织映射（infinitesimal intertwining map）。这些空间就是 $L^2(\mathrm{Whit})^\pm$，即正、负 Whittaker 模型（Whittaker model）。由此可分析 $\omega_{\operatorname{SL}_2}$ 在 $L^2(\mathrm{Whit})^\pm$ 中的谱，并确定哪些 $\lambda$ 使得 $L^2(\mathrm{Whit})_\lambda^\pm\ne0$。可见，$\widetilde{\operatorname{SL}}_2$ 在 $L^2(\mathrm{Whit})^+$ 中的离散谱与其在 $L^2(\mathrm{Whit})^-$ 中的离散谱不相交。因此，$F_Q^+(\lambda)$ 和 $F_Q^-(\lambda)$ 是 $F_Q(\lambda)$ 的稳定且不等价的 $\widetilde{\operatorname{SL}}_2\times O(Q)$ 子模，见上述 (b)。具体来说，这些特征空间中的函数 $f$ 具有如下性质：$f$ 和其 Fourier 变换 $\widehat f$ 的支集都完全位于光锥（light cone）的同一侧。

下一步考察 $\varphi$ 的角向部分 $\rho_2^\pm$。$\rho_2^\pm$ 只是 $L^2(\text{双曲面})$ 上的 Laplace--Beltrami 算子 $W_\xi^\pm$ 的 $L^2$ 特征函数。利用上述 $\pi_Q(\omega_{O(Q)})$ 与 $\pi_Q(\omega_{\operatorname{SL}_2})$ 的关系，以及“双曲面”分析的各种方法 [21]、[3]、[12] 中的任一种，定理2.1中剩余的不可约性问题归结为 $O(Q)$ 在 $L^2(\text{双曲面})$ 中的离散谱的类似问题。

上述讨论还可以确定 $F_Q(s^2-2s)$ 中 $\widetilde K\times K$ 有限函数 $\varphi$ 的正则性与增长性质。特别地，任何这样的函数都唯一地延拓为 $\R^k-\{0\}$ 上的连续函数。剩下需要考察的是原点处的行为。若 $s>\tfrac12k$，则这些函数在原点也连续。这些结论可由 $\varphi$ 的显式表达式推出。类似地，对于 $F_Q^+(s^2-2s)$ 中的 $\widetilde K\times K$ 有限函数 $\varphi$，可以推出下述增长性质：\footnote{\textbf{译者注（本次校订）：}式 (2-1) 的指数中应为 $r_G^{-2}Q(Z,Z)$，并须在右端乘上依赖于 $\varphi$ 的常数 $C_\varphi$。事实上，$\pi_Q(a(r)^{-1})\varphi(Z)=r^{-k/2}\varphi(Z/r)$，故高斯因子的尺度为 $Q(Z,Z)/r^2$；这也与 (3-4) 中 $e^{-c(r_G)^{-2}}$ 的尺度一致。常数 $C_\varphi$ 则不能省略，因为可任意倍乘 $\varphi$。}
\begin{equation}
\begin{split}
 |\pi_Q((G,g)^{-1})(\varphi)(Z)|
 \le{}&r_G^{s-2}\|g^{-1}\|_k^{s+k/2-2}
 \exp\bigl(-\tfrac12\pi r_G^2Q(Z,Z)\bigr)\\
 &\cdot\left(\frac1{\|Z_+\|}\right)^{s+k/2-2},\qquad Z\in\Omega_+.
\end{split}\tag{2-1}
\end{equation}
这里 $Q(Z)=\|Z_+\|^2-\|Z_-\|^2$，$\|\ \|_k$ 是线性算子 $g\in O(Q)$ 的 Frobenius 范数，$r_G$ 是 $G$ 的 Iwasawa 分解（Iwasawa decomposition）的 $A$ 分量，即 $G=K(G)A(r_G)N(x_G)$。

\origsection{3}{$\widetilde{\operatorname{SL}}_2$ 尖点形式（Cusp forms）}
[22] 中 Weil 表示 $\pi_Q$ 的关键性质如下：给定格 $L\subseteq\R^k$，满足 $Q(L,L)\subseteq\Z$，则存在 $\widetilde{\operatorname{SL}}_2\times O(Q)$ 的一个“算术离散”子群 $\Gamma(Q,L)$，使对于 $\R^k$ 上任意 Schwartz 函数 $\varphi$，
\begin{equation}
 \Theta_{\varphi,\eta}^L(G,g)=\sum_{\xi\in L}
 \pi_Q((G,g)^{-1})(\varphi)(\xi+\eta),\qquad\eta\in L_*/L,\tag{3-1}
\end{equation}
是 $\widetilde{\operatorname{SL}}_2\times O(Q)$ 上的 $C^\infty$ 函数，并在右作用下按照 $\Gamma(Q,L)$ 的一个有限维酉表示变换。注意 (1-4) 是它的特例。这一结论的主要分析内容是：Poisson 分布族
\[
 G_\eta=\sum_{\xi\in L}\delta_{\xi+\eta}
\]
定义了 Schwartz 空间 $S(\R^k)$ 到 $\widetilde{\operatorname{SL}}_2\times O(Q)/\Gamma(Q,L)$ 上某个向量丛（vector bundle）的 $C^\infty$ 截面空间的连续交织映射。

上述方法尚不能直接用于 $F_Q^\pm(s^2-2s)$ 上的 $\widetilde K\times K$ 有限函数，因为这些函数不是 Schwartz 函数。结合前面得到的增长估计、连续性和经典 Poisson 求和公式（Poisson summation formula），则可得到以下定理。

\begin{originaltheorem}{定理 3.1（Theorem 3.1）[13]}
设 $s>\tfrac12k$。那么，对任意 $\varphi\in F_Q^\pm(s^2-2s)$ 中的 $\widetilde K\times K$ 有限函数，
\begin{equation}
 \Theta_\varphi^L(G,g)=\sum_{\xi\in L}\pi_Q((G,g)^{-1})(\varphi)(\xi)\tag{3-2}
\end{equation}
是绝对收敛级数。对任意 $(\Gamma,\gamma)\in\Gamma^L(Q)\times O(Q)_L$，其中
\[
 \Gamma^L(Q)=\left\{\left(\begin{pmatrix}a&b\\c&d\end{pmatrix},\varepsilon\right):
 \substack{a,b,c,d\in\Z,\ ad-bc=1,\\c\equiv0\pmod{2N_L},\ b\equiv0\pmod2}\right\},
\]
有函数方程
\begin{equation}
 \Theta_\varphi^L(G\Gamma,g\gamma)=c(\Gamma,\gamma)\Theta_\varphi^L(G,g),\tag{3-3}
\end{equation}
其中 $c(\Gamma,\gamma)$ 为 $\Gamma^L(Q)\times O(Q)_L$ 上取值于 $S_4=\{z\in\C:z^4=1\}$ 的酉特征。此外，$\Theta_\varphi^L$ 是 $\widetilde{\operatorname{SL}}_2\times O(Q)$ 上的 $C^\infty$ 函数，对其泛包络代数（universal enveloping algebra）中的任意 $D$，均满足
\[
 D*\Theta_\varphi^L(G,g)=\Theta_{\pi_Q(D)\varphi}^L(G,g),
\]
其中 $*$ 表示左微分。
\end{originaltheorem}

特别地，由定理3.1可知，$\Theta_\varphi^L$ 是 $\widetilde{\operatorname{SL}}_2\times O(Q)$ 的泛包络代数的中心 $\mathfrak z(\widetilde{\operatorname{SL}}_2\times O(Q))$ 的特征函数，即
\[
 \omega_{\operatorname{SL}_2}*\Theta_\varphi^L(G,g)=(s^2-2s)\Theta_\varphi^L(G,g).
\]
另一方面，由增长估计 (2-1)，有
\begin{equation}
 |\Theta_\varphi^L(G,g)|\le M\|g^{-1}\|_k^{s+k/2-2}
 \begin{cases}r_G^{s-1/2},&r_G\ge1,\\r_G^s e^{-c(r_G)^{-2}},&r_G<1,\end{cases}\tag{3-4}
\end{equation}
其中 $M,c$ 为与 $(G,g)$ 无关的正常数。由此得到以下推论。
\begin{originaltheorem}{定理 3.1 的推论（Corollary to Theorem 3.1）}
在定理3.1的相同假设下，$\Theta_\varphi^L(G,g)$ 关于 $\widetilde{\operatorname{SL}}_2$ 变量是尖点形式。
\end{originaltheorem}

\paragraph*{例 3-1（Example 3-1）.}
$F_Q^+(s^2-2s)$ 中的一个 $\widetilde K\times K$ 有限函数例子是
\begin{equation}
 \varphi(X)=\begin{cases}
 0,&X\in\Omega_-,\\
 Q(X)^{s-1}e^{-\pi Q(X)}\|X_+\|^{-(s+s_1+s_2+k/2-2)}
 Q(X,\xi_+)^{s_1}Q(X,\xi_-)^{s_2},&X\in\Omega_+,
 \end{cases}\tag{3-5}
\end{equation}
其中 $\xi_+$（分别为 $\xi_-$）是 $\C^a$（分别为 $\C^b$）中的非零复各向同性向量。这里的各向同性是相对于 $\C^a$ 上的 $\|\ \|_+^2$（分别为 $\C^b$ 上的 $\|\ \|_-^2$）而言，$s_1,s_2$ 为满足 $s_1-s_2=s-\tfrac12(a-b)$ 的正整数。\footnote{\textbf{译者注（本次校订）：}所述非零复各向同性向量要求相应空间的维数至少为 $2$。当 $b=1$ 时，$\C^b$ 中不存在这样的 $\xi_-$；要包含后文的这一情形，应允许 $s_2=0$ 并省略 $Q(X,\xi_-)^{s_2}$ 这一因子。}令
\begin{equation}
 \widetilde\Theta_\varphi^L(z,g)
 =\Theta_\varphi^L\!\left(\left(\begin{pmatrix}x&y\\0&x^{-1}\end{pmatrix},1\right),g\right)
 (\operatorname{Im}z)^{s/2},\tag{3-6}
\end{equation}
便可用上半平面 $H$ 的经典坐标来表示 $\Theta_\varphi^L$，其中 $z=i/x^2-y/x\in H$。于是
\begin{equation}
 \widetilde\Theta_\varphi^L(z,g)
 =\sum_{\substack{n\in\Z\\n\ge1}}n^{s-1}e^{\pi inz}\varphi_n^{s_1}(g),\tag{3-7}
\end{equation}
其中
\[
 \varphi_n^{s_1}(g)=\sum_{\substack{M\in L\\Q(M,M)=n}}
 Q(M,g^{-1}\xi_+)^{s_1}Q(M,g^{-1}\xi_-)^{s_2}
 \|(gM)_+\|^{-(s+s_1+s_2+k/2-2)}.
\]
那么 $\widetilde\Theta_\varphi^L(z,g)$ 关于 $z$ 是权为 $s$、乘子（multiplier）为 $\nu_Q$、对应于算术群
\[
 \Delta_{N_L}=\{G\in\operatorname{SL}_2(\Z):c_G\equiv0\pmod{2N_L},\ b_G\equiv0\pmod2\}
\]
的全纯尖点形式，即
\begin{equation}
 \widetilde\Theta_\varphi^L\!\left(\frac{az+b}{cz+d},g\right)
 =\nu_Q(G)(cz+d)^s\widetilde\Theta_\varphi^L(z,g),\tag{3-8}
\end{equation}
其中
\[
 G=\begin{pmatrix}a&b\\c&d\end{pmatrix}\in\Delta_{N_L},\qquad
 \nu_Q(G)=(\overline\varepsilon_{d_G})^k
 \left(\frac{2c_G}{d_G}\right)^k\left(\frac{D_{Q(L)}}{d_G}\right),\quad c_G\ne0,
\]
\[
 \varepsilon_m=\begin{cases}1,&m\equiv1\pmod4,\\i,&m\equiv3\pmod4,\end{cases}
\]
括号中的分式表示广义二次剩余符号；分母为奇合数时，应按 Jacobi 符号理解。上述表达式正是 $\widetilde\Theta_\varphi^L$ 关于 $z$ 在 $\{\infty\}$ 处的 Fourier 展开。因此，第 $n$ 个 Fourier 系数 $n^{s-1}\varphi_n^{s_1}(g)$ 是 $O(Q)$ 上相对于群 $O(Q)_L$ 的自守形式。

\origsection{4}{$\Theta_\varphi^L$ 关于 $O(Q)$ 变量的常数项（Constant term）}
主要问题是确定 $\Theta_\varphi^L(G,g)$ 关于 $O(Q)$ 变量的 Fourier 系数。一般情形可写出类似于 Poincaré 级数 Fourier 系数的形式表达式：它是涉及 Bessel 函数和 Kloosterman 型三角和的无穷和。从算术观点看，这样的表达式并不令人满意。

不过，其第零 Fourier 系数，即沿有理 $O(Q)$ 抛物子群的各种幺幂根基（unipotent radical）的常数项，可用较初等的方法确定。

以 $\Q$ 表示有理数域，令
\[
 O(Q)_{\Q\cdot L}=\{\gamma\in O(Q):\gamma(\Q\cdot L)=\Q\cdot L\},
\]
其中 $\Q\cdot L$ 是 $\{r\xi:r\in\Q,\xi\in L\}$ 的有理线性张成空间，其有理维数为 $k$。

我们知道，$O(Q)$ 的每个极大抛物子群（maximal parabolic subgroup）均在 $O(Q)$ 中共轭于形如 $p_F$ 的群，其中
\[
 p_F=\{g\in O(Q):g(F)=F\},
\]
$F$ 为 $\R^k$ 的非零 $Q$ 各向同性子空间，即对所有 $X,X'\in F$，$Q(X,X')=0$。由 Witt 定理，存在另一个 $Q$ 各向同性子空间 $F'$，使 $\dim F=\dim F'$，$F\cap F'=0$，且 $Q$ 在直和 $F+F'$ 上非退化。$p_F$ 的 Levi 分解（Levi decomposition）为
\[
 p_F=\bigl(\operatorname{GL}(F)\times O(Q,(F+F')^\perp)\bigr)\cdot N_F,
\]
其中 $\operatorname{GL}(F)$ 为 $F$ 上可逆实线性变换组成的群，$O(Q,(F+F')^\perp)$ 为 $Q$ 在 $(F+F')^\perp$ 上限制的正交群，$N_F$ 为 $p_F$ 的幺幂根基。这里 $\operatorname{GL}(F)$ 通过映射
\[
 g(u_1+u_2+u_3)=g(u_1)+(g^t)^{-1}(u_2)+u_3,
 \quad u_1\in F,\ u_2\in F',\ u_3\in(F+F')^\perp
\]
嵌入 $O(Q)$。

如果存在 $F'$，以及分别位于 $F,F',(F+F')^\perp$ 中的整格 $L_F,L_{F'},L_{(F+F')^\perp}$，使得 $L$ 与 $L_F\oplus L_{F'}\oplus L_{(F+F')^\perp}$ 在 $\R^k$ 中可公度（commensurate），则称 $p_F$ 与格 $L$ 相容（compatible）。特别地，相对于 $\R^k$ 上由 $\Q\cdot L$ 给出的有理结构，$p_F\cap O(Q)_L$ 是 $p_F$ 的算术子群。

为了研究 $\Theta_\varphi^L(G,g)$ 的尖点性质，必须考察形如
\begin{equation}
 \int_{N_F/(N_F\cap\gamma O(Q)_L\gamma^{-1})}
 \Theta_\varphi^L(G,gv\gamma)\,d\sigma(v)\tag{4-1}
\end{equation}
的积分，其中 $g\in O(Q)$，$\gamma\in O(Q)_{\Q\cdot L}$，$d\sigma$ 为上述商空间上的某个 $N_F$ 不变测度。假设 $p_F$ 与 $L$ 相容；此时该商是紧的。若对于所有 $g\in O(Q)$、$\gamma\in O(Q)_{\Q\cdot L}$ 及与 $L$ 相容的 $p_F$，这些积分均为零，那么 $\Theta_\varphi^L(G,g)$ 就是 $O(Q)$ 上相对于 $O(Q)_L$ 的尖点形式，见 [5]。

另一方面，取 $\varphi\in F_Q^\pm(s^2-2s)_{\widetilde K\times K}$，并在 $s>\tfrac12k$ 时定义
\begin{equation}
 \Phi_\varphi^F(W)=\int_F\varphi(W+U)\,dU,\tag{4-2}
\end{equation}
其中 $W\in(F+F')^\perp$，$dU$ 为 $F$ 上的一个欧几里得测度。我们知道，仿射平面 $W+F$ 是 $\R^k$ 中非零向量 $W$ 的 $N_F$ 轨道。因此，特别地，$\Phi_\varphi^F$ 可以理解为 $\varphi$ 沿 $W$ 的 $N_F$ 轨道的积分；此轨道带有 $N_F$ 不变测度，可与上述 $dU$ 等同。由此得到：
\begin{originaltheorem}{引理 4.1（Lemma 4.1）[13]}
设 $s>\tfrac12k$。那么映射 $\varphi\mapsto\Phi_\varphi^F$ 是从
$F_Q^\pm(s^2-2s)_{\widetilde K\times K}$ 满射到
$F_{Q_F}^\pm(s^2-2s)_{\widetilde K\times K_F}$ 的无穷小
$\widetilde{\operatorname{SL}}_2\times O(Q,(F+F')^\perp)$ 交织映射，其中
\[
 K_F=K\cap O(Q,(F+F')^\perp),
\]
$Q_F$ 是 $Q$ 在 $(F+F')^\perp$ 上的限制。所谓“无穷小”，是指相对于相应泛包络代数表示而言。
\end{originaltheorem}

\paragraph*{注 4.1（Remark 4.1）.}
若 $Q_F$ 正定，则 $F_{Q_F}^+(s^2-2s)$ 的定义与不定 $Q$ 的 $F_{Q_F}^+(\cdots)$ 的定义相似。

引理4.1尤其适用于 $F_{Q_F}^+(s^2-2s)=\{0\}$ 时。例如，当 $Q_F$ 的符号为 $(m,1)$ 时，便有 $F_{Q_F}^-(s^2-2s)=\{0\}$。于是：
\begin{originaltheorem}{引理 4.1 的推论（Corollary to Lemma 4.1）}
设 $(G,g)\in\widetilde{\operatorname{SL}}_2\times O(Q)$。那么在以下情形，$\Phi_{\pi_Q((G,g)^{-1})(\varphi)}^F\equiv0$：
\begin{enumerate}
\item[(a)] $\varphi\in F_Q^-(s^2-2s)_{\widetilde K\times K}$，$\dim F=b-1$；
\item[(b)] $\varphi\in F_Q^-(s^2-2s)_{\widetilde K\times K}$，$\dim F=b$；
\item[(c)] $\varphi\in F_Q^+(s^2-2s)_{\widetilde K\times K}$，$\dim F=b$，且 $b=a-1$ 或 $b=a$；
\item[(d)] $\varphi\in F_Q^+(s^2-2s)_{\widetilde K\times K}$，$\dim F=b-1$，且 $b=a$ 或 $a=2$。
\end{enumerate}
\end{originaltheorem}

需要注意，推论允许 $g\in O(Q)$，而引理4.1中的映射 $\varphi\mapsto\Phi_\varphi^F$ 只对 $O(Q,(F+F')^\perp)$ 交织。原因在于，对于所有 $W\in(F+F')^\perp$ 和所有 $\varphi\in F_Q^+(s^2-2s)_{\widetilde K\times K}$，$\Phi_\varphi^F(W)\equiv0$ 蕴含
\[
 \Phi_{\pi_Q(g)^{-1}\varphi}^F(W)
 =\Phi_{\pi_Q(k)^{-1}(\varphi)}^F\bigl((pW)_{(F+F')^\perp}\bigr)\equiv0.
\]
这里使用相对于抛物子群的分解 $g=k\cdot p$，其中 $k\in K$、$p\in p_F$；$(pW)_{(F+F')^\perp}$ 为 $pW$ 在 $(F+F')^\perp$ 上的投影。

由此可确定 $\Theta_\varphi^L$ 的尖点性质。
\begin{originaltheorem}{定理 4.2（Theorem 4.2）[13]}
令 $p_F$ 为与 $L$ 相容的 $O(Q)$ 的极大抛物子群。取
$\varphi\in F_Q^\pm(s^2-2s)_{\widetilde K\times K}$，其中 $s>\tfrac12k$；取
$\gamma\in O(Q)_{\Q\cdot L}$。令 $S_{\gamma,F}(L)$ 为格
$\gamma(L)\cap(F+(F+F')^\perp)$ 相对于有理分解
$\R^k=F+F'+(F+F')^\perp$ 投影到 $(F+F')^\perp$ 所得的格。那么，对任意 $(G,g)\in\widetilde{\operatorname{SL}}_2\times O(Q)$，
\begin{equation}
 \int_{N_F/(N_F\cap\gamma O(Q)_L\gamma^{-1})}
 \Theta_\varphi^L(G,gv\gamma)\,d\sigma(v)
 =c_1\sum_{\xi\in S_{\gamma,F}(L)}
 \Phi_{\pi_Q((G,g)^{-1})(\varphi)}^F(\xi),\tag{4-3}
\end{equation}
其中 $c_1$ 为仅依赖于 $p_F$ 的正常数。
\end{originaltheorem}

以下交换图示意了 (4-3) 的归纳性质；下方横箭头表示沿 $N_F/(N_F\cap\gamma O(Q)_L\gamma^{-1})$ 取常数项：
\begin{equation}
\begin{tikzcd}[column sep=large,row sep=huge]
 \pi_Q((G,g)^{-1})(\varphi)
 \arrow[r,mapsto]\arrow[d,mapsto,"\text{关于 }L\text{ 的 theta 级数}"']
 &\Phi_{\pi_Q((G,g)^{-1})(\varphi)}^F
 \arrow[d,mapsto,"\text{关于 }S_{\gamma,F}(L)\text{ 的 theta 级数}"]\\
 \Theta_\varphi^L(G,g)
 \arrow[r,mapsto,"\text{取常数项}"']
 &\Theta_{\Phi_{\pi_Q(k_1^{-1})\varphi}^F}^{S_{\gamma,F}(L)}(G,g_1)
\end{tikzcd}\tag{4-4}
\end{equation}
图中右上角还满足
\[
 \Phi_{\pi_Q((G,g)^{-1})(\varphi)}^F
 =\pi_{Q_F}((G,g_1)^{-1})\bigl(\Phi_{\pi_Q(k_1^{-1})\varphi}^F\bigr),
\]
其中 $g=k_1g_1n$，$k_1\in K$，$g_1\in O(Q,(F+F')^\perp)$，$n\in N_F$。\footnote{\textbf{译者注（本次校订）：}这里对任意 $g$ 的分解漏掉了 Levi 因子 $\operatorname{GL}(F)$。一般应写成 $g=k_1ag_1n$，其中 $a\in\operatorname{GL}(F)$。在 (4-2) 中对 $U$ 换元，还会产生 $|\det(a|_F)|^{-1}$；因此上式以及前面的 $g=kp$ 变换式都须计入这一因子。这不影响相应积分恒为零的结论。}由 (4-4) 可见，$\Theta_\varphi^L(G,g)$ 沿相应幺幂商空间的常数项，是一个更低维空间上的函数的 theta 级数。

定理4.2的证明思路是：在 (4-1) 中用 (3-2) 替代 $\Theta_\varphi^L$，并交换求和与积分次序，再按格 $L$ 中的 $N_F\cap\gamma O(Q)_L\gamma^{-1}$ 轨道分类，将结果写成 $\varphi$ 的“$N_F$ 极限圈投影”（horocyclic projection）之和。这里极限圈投影指 $\varphi$ 沿 $\R^k$ 中一个带 $N_F$ 不变测度的 $N_F$ 轨道的积分。(4-2) 就给出了一个例子。本质上，$\R^k$ 中有两类 $N_F$ 轨道。第一类对应于“弯曲的圆锥截面”，称为正则极限圈（regular horocycle）。$\varphi$ 对应于正则极限圈的投影为零；这一陈述在 [13] 中称为第一尖点消失定理（First Cusp Vanishing Theorem）。其关键在于：$\widehat\varphi$ 在 $\R^k$ 上连续，并在光锥上为零。第二类 $N_F$ 轨道就是上面讨论的仿射类型；其极限圈投影的贡献由引理4.1确定。

由引理4.1的推论和定理4.2，得到：
\begin{originaltheorem}{定理 4.2 的推论（Corollary to Theorem 4.2）}
设 $\varphi\in F_Q^-(s^2-2s)_{\widetilde K\times K}$，$s>\tfrac12k$。若 $p_F$ 与 $L$ 相容且 $\dim F=b-1$ 或 $b$，则 $\Theta_\varphi^L(G,g)$ 在
$N_F/(N_F\cap\gamma O(Q)_L\gamma^{-1})$ 方向的常数项消失，$\gamma\in O(Q)_{\Q\cdot L}$。特别地，若 $b=2$ 且 $\varphi\in F_Q^-(s^2-2s)_{\widetilde K\times K}$，则 $\Theta_\varphi^L(G,g)$ 是 $O(Q)$ 上相对于 $O(Q)_L$ 的尖点形式。
\end{originaltheorem}

\paragraph*{注 4.2（Remark 4.2）.}
设 $\varphi\in F_Q^+(s^2-2s)_{\widetilde K\times K}$，$s>\tfrac12k$；$p_F$ 与 $L$ 相容。若 $\dim F=b-1$ 且 $a=2$ 或 $a=b$，或者 $\dim F=b$ 且 $b=a-1$ 或 $b=a$，则相应方向的常数项消失。特别地，若 $a=2$、$b=1$ 或 $2$，并且 $\varphi\in F_Q^+(s^2-2s)_{\widetilde K\times K}$，则 $\Theta_\varphi^L(G,g)$ 是 $O(Q)$ 上相对于 $O(Q)_L$ 的尖点形式。

下面着重讨论 $b=2$ 的情形。此时已知 $O(Q)/K_1$ 是 Hermitian 对称空间（Hermitian symmetric space），其中 $K_1=O(a)\times\operatorname{SO}(2)$。令 $\mathfrak g=\mathfrak k+\mathfrak p$ 为 $O(Q)$ 的 Lie 代数的 Cartan 分解（Cartan decomposition）。复化为 $\mathfrak g_\C=\mathfrak g\otimes_\R\C$ 后，有直和
\[
 \mathfrak g_\C=\mathfrak k_\C\oplus\mathfrak p^+\oplus\mathfrak p^-,
\]
其中 $\mathfrak p^+$ 和 $\mathfrak p^-$ 分别为 $O(Q)/K_1$ 的“原点”处的全纯与反全纯切向量。回忆 $O(Q)$ 的一族全纯离散系列表示（holomorphic discrete series representation）的构造。令 $\chi_n:K_1\to S^1$ 为在 $O(a)$ 上平凡、并将
\[
 \operatorname{SO}(2)=\left\{\begin{pmatrix}\cos\theta&\sin\theta\\-\sin\theta&\cos\theta\end{pmatrix}:
 -\pi<\theta\le\pi\right\}
\]
映到 $e^{-in\theta}$ 的酉特征。构造全纯诱导表示空间（holomorphically induced representation space）
\[
\begin{split}
 \mathcal H(O(Q)/K_1,\chi_n)=\{\varphi:O(Q)\to\C:\;&\varphi\in C^\infty(O(Q)),\\
 &\varphi(gk)=\varphi(g)\chi_n(k)\ (g\in O(Q),k\in K_1),\\
 &\varphi*W\equiv0\ (W\in\mathfrak p^+),\\
 &\int_{O(Q)/K_1}|\varphi(g)|^2d\tau(g)<\infty\}.
\end{split}
\]
其中 $*W$ 表示右卷积，$d\tau$ 为 $O(Q)/K_1$ 上的某个 $O(Q)$ 不变测度。于是：
\begin{originaltheorem}{定理 4.3（Theorem 4.3）[13]、[15]}
$O(Q)$ 在 $\mathcal A_s^-$（见注2.1）中的表示，等价于它在
$\mathcal H(O(Q)/K_1,\chi_{s_2})_\infty$ 中的“全纯”诱导表示，其中
$s_2=|s|+\tfrac12k-2$，下标 $\infty$ 表示 $C^\infty$ 向量。
\end{originaltheorem}

\origsection{5}{Shimura--Niwa 对应（The Shimura--Niwa correspondence）}
前面已经确定 $\Theta_\varphi^L$ 关于两个群变量的尖点性质（其中 $\varphi\in F_Q^\pm(s^2-2s)_{\widetilde K\times K}$，$s>\tfrac12k$），现在用 $\Theta_\varphi^L$ 作积分核，继续研究两个群上的模形式之间的对应。特别地，取例3-1中的 $\varphi$，并考虑 $\widetilde\Theta_\varphi^L(z,g)$，见 (3-7)。令 $f$ 为 $H$ 上权为 $s$（$s>\tfrac12k$）的全纯尖点形式，满足
\[
 f(\gamma\cdot z)=\nu_Q(\gamma)(c_\gamma z+d_\gamma)^s f(z),\qquad\gamma\in\Delta_{N_L}.
\]
如第1节，考虑 $f$ 与 $\widetilde\Theta_\varphi^L$ 的 Petersson 内积：
\begin{equation}
 \langle\widetilde\Theta_\varphi^L(\ ,g)\mid f(\ )\rangle
 =\int_{\Delta_{N_L}\backslash H}\widetilde\Theta_\varphi^L(z,g)
 \overline{f(z)}|\operatorname{Im}z|^{s-2}\,dx\,dy.\tag{5-1}
\end{equation}
(5-1) 的定义乍看与 (1-8) 的 Petersson 内积
$\langle\widetilde\Theta_{\psi_*}^L(\ ,g)\mid f(\ )\rangle$ 不同；后一个内积中的 theta 级数由 Schwartz 函数 $P(X)e^{-\pi[X,X]}$ 构造。沿用第1节的记号，当 $\eta\in L$ 时，省略 $\widetilde\Theta_{\psi,\eta}^L$ 的下标 $\eta$。因此，需要弄清 $\widetilde\Theta_{\psi_*}^L$ 与上面定义的 $\widetilde\Theta_\varphi^L$ 之间的关系。不过，映射 $\psi\mapsto\Theta_\psi^L$ 是从 $L^2(\R^k)$ 中某些 $\widetilde{\operatorname{SL}}_2\times O(Q)$ 稳定子空间到该群上的 $C^\infty$ 函数空间的无穷小交织映射。这就引出了一个自然的问题：给定 $\widetilde K\times K$ 有限函数 $\psi$，令 $\psi_D$ 为它在 $L^2(\R^k)$ 中 $\widetilde{\operatorname{SL}}_2\times O(Q)$ 离散谱上的投影，theta 级数 $\Theta_\psi$ 与 $\Theta_{\psi_D}$ 有何关系？以下引理给出部分答案。

\begin{originaltheorem}{引理 5.1（Lemma 5.1）[16]}
设 $F$ 是 $L^2(\R^k)$ 中的 $K$ 有限 Schwartz 函数，并假设
\[
 \pi_m(k(\theta,\varepsilon))(F)=e^{-is'\theta}\varepsilon^kF,
 \quad-\pi<\theta\le\pi,\quad\varepsilon=\pm1.
\]
设 $s'>\tfrac12k+1$。以 $P_{s'}^+$ 表示 $F_Q$ 到子空间
$F_Q^+(s'^2-2s')$ 的投影。那么
\begin{equation}
 \langle\widetilde\Theta_F^L(\ ,g)\mid f(\ )\rangle
 =\langle\widetilde\Theta_{P_{s'}^+(F)}^L(\ ,g)\mid f(\ )\rangle.\tag{5-2}
\end{equation}
\end{originaltheorem}

\paragraph*{注 5.1（Remark 5.1）.}
若 $s'<-(\tfrac12k+1)$，则有类似陈述，其中 $P_{s'}^+$ 换为 $P_{s'}^-$，即 $F_Q$ 到 $F_Q^-(s'^2+2s')$ 的投影。因此，(1-8) 与 (5-1) 这两个对应本质上相同。

下一个主要问题，是刻画映射
\[
 f\longmapsto\langle\widetilde\Theta_\varphi^L(\ ,g)\mid f(\ )\rangle
\]
在 $f$ 遍历空间 $[\Delta_{N_L},s,\nu_Q]_0$ 时的像；该空间由 $H$ 上全纯、对所有 $\gamma\in\Delta_{N_L}$ 满足上述变换律，并在 $\Delta_{N_L}$ 于 $\Q\cup\{\infty\}$ 中的尖点处消失的函数构成，即权为 $s$、乘子为 $\nu_Q$ 的全纯尖点形式。首先，取 Eisenstein--Poincaré 级数（Eisenstein--Poincaré series）
\[
 f(z)=G_n(z)=\sum_{\gamma\in(\Delta_{N_L})_\infty\backslash\Delta_{N_L}}
 \left(\frac1{c_\gamma z+d_\gamma}\right)^s
 \overline{\nu_Q(\gamma)}e^{\pi in\gamma(z)},
\]
则
\[
 \langle\widetilde\Theta_\varphi^L(\ ,g)\mid G_n(\ )\rangle
 =\bar c_1\,\overline{\varphi_n^{s_1}(g)},
\]
其中 $c_1$ 是与 $g,n$ 无关的非零常数。\footnote{\textbf{译者注（本次校订）：}按 (5-1) 的内积约定，内积对第一个变量线性。展开 Poincaré 级数并积分，应得到常数乘以 $\varphi_n^{s_1}(g)$，而不是它的复共轭。相应地，下句所述像空间应由 $\varphi_n^{s_1}$ 张成；若改用相反的内积次序，才会出现所写的复共轭。}由于 $s>\tfrac12k$，函数 $G_n$ 张成 $[\Delta_{N_L},s,\nu_Q]_0$，所以这些内积构成的空间恰为 $\overline{\varphi_n^{s_1}}$（$n\ge1$）的复线性张成空间。

由定理4.2的推论，可知哪些情形下所有 $\varphi_n^{s_1}$ 都是 $O(Q)$ 上相对于 $O(Q)_L$ 的尖点形式。$\varphi_n^{s_1}$ 一般 Fourier 系数具有 Poincaré 级数式的复杂结构，难以用算术数据直接描述。不过，利用第4节的结果，可以确定哪些 $f\in[\Delta_{N_L},s,\nu_Q]_0$ 的上述内积是 $O(Q)$ 上相对于 $O(Q)_L$ 的尖点形式。

取极大 $Q$ 各向同性子空间 $F_j$，使 $p_{F_j}$ 与 $L$ 相容。考虑 $Q$ 各向同性子空间的旗（flag）
\[
 F_1\subseteq F_2\subseteq\cdots\subseteq F_j,
\]
使每个 $p_{F_i}$ 与 $L$ 相容，且 $F_{i+1}/F_i$ 的维数为 $1$。此时 $O(Q)$ 的有理点有如下分解：
\[
 O(Q)_{\Q\cdot L}=O(Q)_L\cdot\Xi_{F_i}\cdot(p_{F_i})_{\Q\cdot L},
\]
其中 $\Xi_{F_i}$ 为 $O(Q)_{\Q\cdot L}$ 的某个有限子集（见 [2，第104页]）。于是得到尖点形式的下述判据。
\begin{originaltheorem}{定理 5.2（Theorem 5.2）}
设 $s>\tfrac12k$，$\varphi$ 由 (3-5) 给定，$f\in[\Delta_{N_L},s,\nu_Q]_0$。令 $Y_{F_i}$ 为
$F_{Q_{F_i}}^+(s^2-2s)$ 中唯一的
$K_{F_i}=K\cap O(Q,(F_i+F_i')^\perp)$ 不变子空间，由
$\{\Phi_{\pi_Q(k)\varphi}^{F_i}:k\in K\}$ 的复线性张成空间给定。那么
$\langle\widetilde\Theta_\varphi^L(\ ,g)\mid f(\ )\rangle$ 是 $O(Q)$ 上相对于 $O(Q)_L$ 的尖点形式，当且仅当
\begin{equation}
 \langle\widetilde\Theta_\psi^{S_{\lambda_{ji},F_i}(L)}(\ ,g_1)\mid f(\ )\rangle\equiv0\tag{5-3}
\end{equation}
对所有 $i=1,\ldots,j$、所有 $\lambda_{ji}\in\Xi_{F_i}^{-1}$、$\psi\in Y_{F_i}$ 和
$g_1\in O(Q,(F_i+F_i')^\perp)$ 成立。
\end{originaltheorem}

注意，$Y_{F_i}$ 有限维，$\Xi_{F_i}$ 是有限集；当 $g$ 遍历 $O(Q,(F_i+F_i')^\perp)$ 时，函数
$\widetilde\Theta_\psi^{S_{\lambda_{ji},F_i}(L)}(z,g)$ 张成 $[\Delta_{N_L},s,\nu_Q]_0$ 的一个有限维子空间。特别地，如果 $Q$ 的符号为 $(k-1,1)$，则只须对 $O(Q,(F+F')^\perp)$ 的单位元 $g_1$ 检验 (5-3)，因为此时该群等于 $K_F$。这也说明了 [1] 中从 $\widetilde{\operatorname{SL}}_2$ 尖点形式提升到 $O(3,1)$ 尖点形式所需的正交性条件。因此，上述内积是 $O(Q)$ 上的尖点形式，当且仅当 $f$ 属于由低维二次型与格给出的有限族对应
\[
 f\longmapsto\langle\widetilde\Theta_\psi^{S_{\lambda_{ji},F_i}(L)}(\ ,g_i)\mid f(\ )\rangle
\]
的共同核。

\origsection{6}{“Zagier 极化恒等式”（Zagier polarization identity）}
回忆，当 $b=2$ 时，$O(Q)/K_1$ 是 Hermitian 对称空间。事实上，可将它实现为管状域（tube domain）。令 $\R^{k-2}$ 为 $e_2,e_3,\ldots,e_{k-2},e_k$ 的张成空间（见第1节）。于是 $Q$ 在 $\R^{k-2}$ 上的限制为 Lorentz 型，即符号为 $(k-3,1)$。令
\[
 \mathcal C_+=\{W\in\R^{k-2}:Q(W,W)<0,\ Q(W,e_k)<0\},
\]
这就是未来光锥，也称前向光锥（future or forward light cone）。其管状域为
\[
 T(\mathcal C_+)=\{X+iY:X,Y\in\R^{k-2},\ Y\in\mathcal C_+\}
 =\R^{k-2}+i\mathcal C_+.
\]
令 $O(Q)^\#$ 为 $O(Q)$ 中对应于 $T(\mathcal C_+)$ 的全体解析自同构群的连通分支的子群。若 $F_1$ 表示 $Q$ 各向同性直线
\[
 \{t(e_1-e_{k-1})/\sqrt2:t\in\R\},
\]
则 $p_{F_1}^\#=p_{F_1}\cap O(Q)^\#$ 在 $T(\mathcal C_+)$ 上如下作用：若
\[
 q=m_1A(r)n_1\in p_{F_1}^\#,
 \quad m_1\in O(Q,(F_1+F_1')^\perp),\ A(r)\in S_{F_1},\ n_1\in N_{F_1},
\]
则 $q(Z)=r\{m_1(Z+Y)\}$，其中 $n_1$ 由一个元素
$Y\in\R^{k-2}=(F_1+F_1')^\perp$ 唯一确定。

回忆，$T(\mathcal C_+)$ 上存在唯一的全纯自守因子（holomorphic automorphy factor）$\mathcal D$，使如上给定的 $m_1A(r)n_1$ 满足
$\mathcal D(m_1A(r)n_1,Z)=1/r$。

由定理4.3，取 $\varphi\in\mathcal A_s^-$（$b=2$ 且 $s>\tfrac12k$）为 $O(Q)$ 相对于 $K_1$ 的“最低权”（lowest weight）向量；除去非零标量倍数后，该向量唯一。其显式形式为
\begin{equation}
 \varphi(X)=\begin{cases}
 0,&X\in\Omega_+,\\
 Q(X,X)^{s-1}e^{\pi Q(X)}Q(X,\xi_+)^{-s_2},&X\in\Omega_-,
 \end{cases}\tag{6-1}
\end{equation}
其中 $\xi_+=e_{k-1}+ie_k$，$s_2=s+\tfrac12k-2$。

接下来对由 $\varphi$ 构造的核作一个扭转，以此得到 (5-1) 的变式。假设 $Q$ 整格 $L$ 有 $Q$ 正交直和分解
\[
 L=\mathcal L\oplus(\Z v\oplus\Z\widetilde v),
\]
其中 $\mathcal L\subseteq\R^{k-2}$（即 $\{e_2,e_3,\ldots,e_{k-2},e_k\}$ 的实张成空间），满足
\[
 Q(\mathcal L)\subseteq2\Z,\qquad n_{\mathcal L}Q(\mu,\mu)\in2\Z
 \quad(\mu\in\mathcal L_*(Q)),
\]
$\mathcal L_*(Q)$ 是 $\mathcal L$ 的 $Q$ 对偶。此外，$v,\widetilde v$ 张成双曲平面，且
$Q(v,v)=Q(\widetilde v,\widetilde v)=0$、$Q(v,\widetilde v)=1$。再假设整数 $t$ 满足
$n_{\mathcal L}\mid t$ 且 $4\mid t$；于是 $L$ 的指数 $n_L$ 等于 $t$。\footnote{\textbf{译者注（本次校订）：}按这里写出的正交直和与 $Q(v,\widetilde v)=1$，双曲格 $\Z v\oplus\Z\widetilde v$ 是幺模格，因此实际有 $n_L=n_{\mathcal L}$，并不能由 $n_{\mathcal L}\mid t$ 推出 $n_L=t$。而且 $v\in L$，故随后所有 $\Theta_{\varphi,rv}^L$ 都等于 $\Theta_\varphi^L$，非平凡 $\chi$ 的加权和便恒为零。要得到本节所讨论的非平凡扭转，格的定义与平移向量必须作相应调整，使这些向量代表不同的对偶格陪集。}

令 $\chi$ 为模 $t$ 的 Dirichlet 特征（Dirichlet character），考虑形式和
\[
 \Theta_{\varphi,\chi}^L(G,g)=\sum_{r\bmod t}\chi(r)\Theta_{\varphi,rv}^L(G,g).
\]
仿照例3-1，定义
\begin{equation}
\begin{split}
 \widetilde\Theta_{\varphi,\chi}^L(z,g)
 &=\Theta_{\varphi,\chi}^L\!\left(\left(\begin{pmatrix}x&y\\0&x^{-1}\end{pmatrix},1\right),g\right)
 (\operatorname{Im}z)^{-|s|/2}\\
 &=\sum_{r\le-1}\beta_{r,\chi}(g)|r|^{s-1}e^{\pi irz},
\end{split}\tag{6-2}
\end{equation}
其中 $z=-y/x-ix^2\in\overline H$，$\overline H$ 为下半平面，而\footnote{\textbf{译者注（本次校订）：}由这里的上三角矩阵及第2节的作用公式可知，下半平面坐标应为 $z=-y/x-i/x^2$。这与 (3-6) 的上半平面坐标互为复共轭，并使 $e^{\pi irz}$ 的衰减因子与 $\pi_Q(G^{-1})$ 中的高斯因子一致。}
\[
 \beta_{r,\chi}(g)=\sum_{u\bmod t}\sum_{\substack{M\in L\\Q(M+uv,M+uv)=r}}
 \chi(u)Q(M+uv,g^{-1}\xi_+)^{-s_2}.
\]
利用全纯自守因子 $\mathcal D$，置
\[
 \widetilde\beta_{r,\chi}(Z)=[\mathcal D(g,ie_k)]^{s_2}\beta_{r,\chi}(g^{-1}),
 \qquad Z=g\,i(\sqrt2e_k)\in T(\mathcal C_+).
\]
由此得到
\begin{equation}
\begin{split}
 \widetilde\beta_{r,\chi}(Z)
 =\sum_{u\bmod t}\sum_{\substack{M\in L\\Q(M+uv,M+uv)=r}}\chi(u)
 \Bigl\{&-\tfrac12Q(M,v/\sqrt2)Q(Z,Z)\\
 &+Q(M,Z)+Q(M,\sqrt2\widetilde v)+\sqrt2u\Bigr\}^{-s_2}.
\end{split}\tag{6-3}
\end{equation}
此外，$\widetilde\beta_{r,\chi}$ 是 $T(\mathcal C_+)$ 上的全纯函数，对所有 $Z\in T(\mathcal C_+)$ 和
\[
 \gamma\in O(Q)_{L,v}^\#=\{\gamma\in O(Q)_L^\#:\gamma(v)\equiv v\pmod L\},
\]
满足
\[
 \widetilde\beta_{r,\chi}(\gamma Z)
 =[\mathcal D(\gamma,Z)]^{s_2}\widetilde\beta_{r,\chi}(Z).
\]
另一方面，$\widetilde\Theta_{\varphi,\chi}^L$ 关于 $z\in\overline H$ 是全纯模尖点形式，即对所有
\[
 \gamma\in\Gamma_0(t)=\left\{\begin{pmatrix}a&b\\c&d\end{pmatrix}\in\operatorname{SL}_2(\Z):c\equiv0\pmod t\right\},
\]
有
\[
 \widetilde\Theta_{\varphi,\chi}^L(\gamma z,g)
 =j(\gamma,z)^{2|s|}\lambda_{Q,|s|}^L(d_\gamma)\chi(d_\gamma)
 \widetilde\Theta_{\varphi,\chi}^L(z,g),
\]
其中
\[
 j(g,z)=\left(\frac{c_g}{d_g}\right)\overline\varepsilon_{d_g}(c_gz+d_g)^{1/2},
 \qquad
 \lambda_{Q,|s|}^L(m)=\left(\frac2m\right)^k
 \left(\frac{D_{Q(L)}}m\right)\left(\frac{-1}m\right)^{s-k/2}
\]
是模 $t$ 的 Dirichlet 特征。

如 (5-1) 定义 Petersson 内积
$\langle\widetilde\Theta_{\varphi,\chi}^L(\ ,g)\mid h_f(\ )\rangle$，其中
$h_f(z)=\overline{f(\bar z)}$，且 $f\in S_{2|s|}(\Gamma_0(t),\beta)$。此空间由满足下列条件的 $\phi:H\to\C$ 构成：(i) 在 $H$ 上全纯；(ii) 对所有 $\gamma\in\Gamma_0(t)$ 和 $z\in H$，
$\phi(\gamma z)=j(\gamma,z)^{2|s|}\beta(d_\gamma)\phi(z)$；(iii) 在 $\Gamma_0(t)$ 于 $\Q\cup\{\infty\}$ 中的各尖点处消失。这里
\[
 \beta(m)=(\lambda_{Q,|s|}^L\otimes\bar\chi)(m)(-1/m)^{2|s|}.
\]
定义
\begin{equation}
 F_f(Z\mid\varphi,L,v,\chi,\Gamma_0(t))
 =\langle\widetilde\Theta_{\varphi,\chi}^L(z,g^{-1})\mid h_f(z)\rangle
 [\mathcal D(g,ie_k)]^{s_2},\tag{6-4}
\end{equation}
其中 $g(i\sqrt2e_k)=Z\in T(\mathcal C_+)$，$s_2=s+\tfrac12k-2>k-2$。于是对所有 $Z\in T(\mathcal C_+)$ 和 $\gamma\in O(Q)_{L,v}^\#$，有
\[
 F_f(\gamma Z\mid\varphi,L,v,\chi,\Gamma_0(t))
 =\mathcal D(\gamma,Z)^{s_2}F_f(Z\mid\varphi,L,v,\chi,\Gamma_0(t)).
\]

现在的主要问题是刻画这个映射的像，关键在于计算 $\widetilde\beta_{r,\chi}$ 的 Fourier 系数。如 [23]、[11]、[13] 所示，这些系数含有复杂的超越函数，是涉及 Bessel 函数与 Kloosterman 型三角和的无穷和。

已知 $N_{F_1}\simeq\R^{k-2}$ 通过平移作用于 $T(\mathcal C_+)$。因此，存在格
$\mathcal L_{F_1}\subset\R^{k-2}=(F_1+F_1')^\perp$，使 $\xi\in\mathcal L_{F_1}$ 时
\[
 F_f(Z+\xi\mid\varphi,L,v,\chi,\Gamma_0(t))
 =F_f(Z\mid\varphi,L,v,\chi,\Gamma_0(t)).
\]
简单计算表明
\[
 \mathcal L_{F_1}=\sqrt2\{\xi\in\mathcal L:Q(\xi,\xi)\equiv0\pmod{2t},\
 Q(\mathcal L,\xi)\equiv0\pmod t\}.
\]
因此有展开
\begin{equation}
 F_f(Z\mid\varphi,L,v,\chi,\Gamma_0(t))
 =\sum_{\mu\in(\mathcal L_{F_1})_*(Q)\cap\mathcal C_+}
 a(\mu,f)e^{-2\pi iQ(\mu,Z)},\tag{6-5}
\end{equation}
其中
\[
 (\mathcal L_{F_1})_*(Q)=\{\mu\in(F_1+F_1')^\perp:Q(\mu,\mathcal L_{F_1})\subseteq\Z\}
\]
是 $\mathcal L_{F_1}$ 的 $Q$ 对偶，并且
\begin{equation}
 a(\mu,f)=\int_{\R^{k-2}/\mathcal L_{F_1}}
 F_f(X+iY\mid\varphi,L,v,\chi,\Gamma_0(t))
 e^{2\pi iQ(X+iY,\mu)}\,dX.\tag{6-6}
\end{equation}
因此，问题归结为利用 $f$ 在 $\{\infty\}$ 处的 Fourier 系数来确定 $a(\mu,f)$。由定理4.2的推论立即知道，若
$\mu\notin(\mathcal L_{F_1})_*(Q)\cap\mathcal C_+$，则 $a(\mu,f)=0$。

先说明这一计算背后的思路。已知空间 $S_{2s}(\Gamma_0(t),\beta)$ 有一个再生核（reproducing kernel）$K:H\times H\to\C$，关于第一变量全纯、第二变量反全纯，对所有 $z_1,z_2\in H$ 和 $\gamma_1,\gamma_2\in\Gamma_0(t)$ 满足
\[
 K(\gamma_1z_1,\gamma_2z_2)
 =j(\gamma_1,z_1)^{2s}\overline{j(\gamma_2,z_2)^{2s}}
 \beta(d_{\gamma_1})\overline{\beta(d_{\gamma_2})}K(z_1,z_2).
\]
此外，对所有 $f\in S_{2s}(\Gamma_0(t),\beta)$，有 $\langle f(z),K(z,w)\rangle=f(w)$，其中内积为 Petersson 内积。当 $s$ 充分大时，有收敛展开
\[
 K(z,w)=\sum_{n\ge1}n^{s-1}e^{-2\pi in\bar w}G_n(z),
\]
其中 $G_n$ 为属于 $S_{2s}(\Gamma_0(t),\beta)$ 的 Eisenstein--Poincaré 级数。另一方面，$\overline{K(z,w)}=K(w,z)$，所以另有展开
\[
 K(z,w)=\sum_{n\ge1}n^{s-1}\overline{G_n(w)}e^{2\pi inz}.
\]
第二个展开可解释为 $K(z,w)$ 关于 $z$ 在 $\{\infty\}$ 处的 Fourier 展开，即 $z\mapsto K(z,w)$ 是 $S_{2s}(\Gamma_0(t),\beta)$ 的一个元素，而第 $n$ 个 Fourier 系数 $n^{s-1}\overline{G_n(w)}$ 是反全纯型模形式。

注意，第二个展开与 (3-6) 给出的 $\widetilde\Theta_\varphi^L(z,g)$ 关于 $z$ 的 Fourier 展开相似，其中每个 Fourier 系数是 $O(Q)$ 上的模形式。这提示我们寻找形如
\[
 \widetilde\Theta_\varphi^L(z,g)=\sum_{n\ge1}n^{s-1}(\varphi_n^{s_1})^*(g)G_n(z)
\]
的“极化恒等式”。这里 $(\varphi_n^{s_1})^*:O(Q)\to\C$ 应满足：对所有 $g\in O(Q)$ 和 $\gamma\in O(Q)_L\cap p_{F_1}$，
\[
 (\varphi_n^{s_1})^*(g\gamma)=(\varphi_n^{s_1})^*(g),
\]
而且 $\varphi_n^{s_1}$ 应由它对 $O(Q)_L/(O(Q)_L\cap p_{F_1})$ 平均得到，即
\[
 \varphi_n^{s_1}(g)=\sum_{\gamma\in O(Q)_L/(O(Q)_L\cap p_{F_1})}
 (\varphi_n^{s_1})^*(g\gamma).
\]
Zagier [23] 在 $a=b=2$ 的情形证明了这样的恒等式。这一公式还可以推广：Oda [11] 处理了 $k\ge6$，[16] 则给出了不要求 $b=2$ 的一般版本。

\begin{originaltheorem}{定理 6.1（Zagier 恒等式，Zagier identity）}
设 $s>2k$，令
\[
 \Xi_{\varphi,\chi}^L(z,Z)=\widetilde\Theta_{\varphi,\chi}^L(z,g^{-1})
 [\mathcal D(g,ie_k)]^{s_2},\qquad Z=g\,i(\sqrt2e_k)\in T(\mathcal C_+),
\]
见 (6-2)。则有公式
\begin{equation}
 \Xi_{\varphi,\chi}^L(z,Z)=\sum_{\substack{n\in\Z\\n\le-1}}
 |n|^{s-1}\beta_{n,\chi}^*(Z)G_n(z,\chi),\tag{6-7}
\end{equation}
其中 $G_n$ 为下半平面 $\overline H$ 上的 Eisenstein--Poincaré 级数：
\begin{equation}
 G_n(z,\chi)=\sum_{\gamma\in\Gamma_0(t)_\infty\backslash\Gamma_0(t)}
 \left(\frac1{c_\gamma z+d_\gamma}\right)^s
 \overline{s_{Q,\chi}^L(\gamma)}e^{2\pi in\gamma(z)},\tag{6-8}
\end{equation}
其中
\[
 s_{Q,\chi}^L(\gamma)=(\overline\varepsilon_{d_\gamma})^k
 \left(\frac{2c_\gamma}{d_\gamma}\right)^k
 \left(\frac{D_{Q(L)}}{d_\gamma}\right)\chi(d_\gamma).
\]
注意，对 $\gamma\in\Gamma_0(t)$，有
\[
 \widetilde\Theta_{\varphi,\chi}^L(\gamma z,g)
 =s_{Q,\chi}^L(\gamma)(c_\gamma z+d_\gamma)^s
 \widetilde\Theta_{\varphi,\chi}^L(z,g).
\]
而
\begin{equation}
 \beta_{n,\chi}^*(Z)=\sum_{v\bmod t}\chi(v)
 \left\{\sum_{\substack{\xi\in\mathcal L,\ Q(\xi,\xi)=n\\j\in\Z}}
 \{Q(\xi,Z)+\sqrt2[jt-v]\}^{-s_2}\right\}.\tag{6-9}
\end{equation}
\end{originaltheorem}

\paragraph*{注 6.1（Remark 6.1）.}
函数 $\beta_{n,\chi}^*$ 由 $\beta_{n,\chi}$ 把 (6-3) 的求和限制到形如
\[
 \{\xi+tjv:\xi\in\mathcal L,\ Q(\xi,\xi)=n,\ j\in\Z\}
\]
的格点而得。这些奇异极限圈格点（singular horocyclic lattice points）的集合在 $O(Q)_L\cap p_{F_1}$ 下稳定。

利用定理6.1可求出 (6-6) 中的 $a(\mu,f)$。做法是对 (6-9) 右端的内层和应用 Lipschitz 恒等式（Lipschitz identity），即\footnote{\textbf{译者注（本次校订）：}这里 Lipschitz 恒等式右端的系数应为 $(-2\pi i)^{s_2}/(s_2-1)!$。当整数 $s_2$ 为奇数时，所列公式差一个负号；从 $s_2=2$ 的恒等式逐次求导即可核对这一符号。}
\[
 \sum_{m\in\Z}\left(\frac1{w+m}\right)^{s_2}
 =\frac{(2\pi i)^{s_2}}{(s_2-1)!}\sum_{r=1}^{\infty}r^{s_2-1}e^{2\pi iwr},\qquad w\in H.
\]
由此推出：
\begin{originaltheorem}{定理 6.1 的推论（Corollary to Theorem 6.1）}
设 $s>2k$，$f\in S_{2s}(\Gamma_0(t),\sigma)$，其中 $\sigma$ 为模 $t$ 的 Dirichlet 特征。令
\[
 f(z)=\sum_{n\ge1}a_f(n)e^{2\pi inz}
\]
为 $f$ 在 $\{\infty\}$ 处的 Fourier 展开。若
$\mu\in(\mathcal L_{F_1})_*(Q)\cap\mathcal C_+$，且
$\mu\notin\mathcal L/(\sqrt2\,t)$，则 $a(\mu,f)=0$。
另一方面，若 $\mu\in(\mathcal L_{F_1})_*(Q)\cap\mathcal C_+$，且
$\mu=(m/(\sqrt2\,t))\xi$，其中 $\xi$ 是格 $\mathcal L$ 的本原元素（primitive element），$m$ 为正整数，则
\begin{equation}
 a(\mu,f)=c_1'G(\chi_\sigma,t)t^{-s_2}
 \sum_{\substack{\nu>0\\\nu\mid m}}\overline{\chi_\sigma(\nu)}\nu^{s_2-1}
 a_f\!\left(\frac{m^2}{\nu^2}|Q(\xi,\xi)|\right),\tag{6-10}
\end{equation}
其中 Dirichlet 特征 $\chi_\sigma$ 给定为
\[
 \chi_\sigma(x)=\overline{\sigma(x)}(-1/x)^{2s}\lambda_{Q,s}^L(x),
\]
Gauss 和（Gauss sum）给定为
\[
 G(\chi_\sigma,t)=\sum_{\nu\bmod t}\chi_\sigma(\nu)e^{2\pi i\nu/t},
\]
$c_1'$ 为仅依赖于 $s$ 的非零常数。\footnote{\textbf{译者注（本次校订）：}将广义 Gauss 和写成 $\overline{\chi_\sigma(r)}G(\chi_\sigma,t)$，并对所有整数 $r$ 使用，要求 $\chi_\sigma$ 为本原特征。对于一般非本原特征，须保留 $\sum_{u\bmod t}\chi_\sigma(u)e^{2\pi iru/t}$，式 (6-10) 不能直接照用。例如模 $4$ 的主特征在 $r=1$ 时的 Gauss 和为 $0$，在 $r=2$ 时却为 $-2$。}
\end{originaltheorem}

沿用自守形式理论中熟知的方法，可对自守函数
$F_f(\ \mid\varphi,L,v,\chi_\sigma,\Gamma_0(t))$ 构造一个 Dirichlet 级数。$k=3,4$ 的情形已在 [10]、[6] 中得到广泛研究，因此以下假设 $k\ge5$。特别地，令
\begin{equation}
 R(\mathfrak s,f)=\sum_{\Omega\in\{(\mathcal L_{F_1})_*(Q)\cap\mathcal C_+\}}
 a(\Omega,f)\frac1{\varepsilon(\Omega)}|Q(\Omega,\Omega)|^{-\mathfrak s},\tag{6-11}
\end{equation}
其中花括号表示 $O(Q,F_1+F_1')^\#\cap O(Q)_L$ 在
$(\mathcal L_{F_1})_*(Q)\cap\mathcal C_+$ 中的等价类集合，$\varepsilon(\Omega)$ 为该群中固定 $\Omega$ 的有限子群的阶。

由 (6-10) 立即得出：
\begin{originaltheorem}{命题 6.2（Proposition 6.2）[16]}
在定理6.1的相同假设下，另设 $k\ge5$。取 $\mathfrak s\in\C$，使 $\operatorname{Re}\mathfrak s$ 充分大。那么有恒等式
\begin{equation}
\begin{split}
 R(\mathfrak s,f)={}&d_1G(\chi_\sigma,t)t^{2\mathfrak s-s_2}2^{\mathfrak s}
 L(\bar\chi_\sigma,2\mathfrak s+1-s_2)\\
 &\cdot\sum_{\substack{n\in\Z\\n\le-1}}
 a_f(-n)M(Q_1,\mathcal L,n)|n|^{-\mathfrak s},
\end{split}\tag{6-12}
\end{equation}
其中 $M(Q_1,\mathcal L,n)$ 是二次型 $Q_1$（即 $Q$ 在 $F_1+F_1'$ 上的限制）相对于格 $\mathcal L$ 在水平 $n$ 二次曲面上的 Siegel 质量数（Siegel mass number），即
\[
 M(Q_1,\mathcal L,n)=\sum_{i=1}^{h(n)}\varepsilon(\xi_i)^{-1},
\]
$\xi_1,\ldots,\xi_{h(n)}$ 遍历
$O(Q,F_1+F_1')^\#\cap O(Q)_L$ 在
$\mathcal L\cap\{X\in\R^{k-2}:Q(X,X)=n\}$ 中的各轨道的一组代表，$d_1$ 是仅依赖 $s_2$ 的非零常数（此处原文回忆 $s_2=s+\tfrac12k-1$）。此外，$L(\delta,\mathfrak s)$ 为与 Dirichlet 特征 $\delta$ 相联系的经典 $L$ 函数。
\end{originaltheorem}
\noindent\textit{译注：本命题中 $Q_1$ 的限制空间在原文印作 $F_1+F_1'$，但相邻公式使用的是 $(F_1+F_1')^\perp$ 上的格；原文此处的 $s_2$ 也与前文 $s_2=s+\tfrac12k-2$ 不同。以上保留原文，供读者辨别。}

这样，$R(\mathfrak s,f)$ 可写成初等函数（形如 $a^{\mathfrak s}$）、一个 $L$ 函数以及两个 Dirichlet 级数的 Rankin 卷积（Rankin convolution）的乘积。这两个级数是
\[
 D(\mathfrak s,f)=\sum_{n\ge1}a_f(n)n^{-\mathfrak s}
\]
和 Siegel zeta 函数
\[
 \zeta_-(Q_1,\mathcal L,\mathfrak s)=\sum_{\substack{n\in\Z\\n\le-1}}
 M(Q_1,\mathcal L,n)|n|^{-\mathfrak s}.
\]
函数 $R(\mathfrak s,f)$ 的解析性质可由上述结果确定。若令
\[
 R^*(\mathfrak s,f)=\pi^{-2\mathfrak s}\Gamma(\mathfrak s-\tfrac12k+2)
 \Gamma(\mathfrak s)R(\mathfrak s,f),
\]
则 $R^*(\mathfrak s,f)$ 可解析延拓到整个 $\mathfrak s$ 平面，$\Gamma$ 表示 Gamma 函数。

\paragraph*{注 6-2（Remark 6-2）.}
利用 (6-12)，可推出 $R(\mathfrak s,f)$ 的某种 Euler 乘积展开（Euler product expansion）。为此，只须研究 $D(\mathfrak s,f)$ 与 $\zeta_-(Q_1,\mathcal L,\mathfrak s)$ 的 Rankin 卷积可能具有的 Euler 乘积性质。若 $k$ 偶，且这两个级数都具有 $\operatorname{GL}_2$ 理论中通常的 Euler 乘积，则其 Rankin 卷积的 Euler 乘积在几乎所有素数 $p$ 处都有分子二次、分母四次的局部因子。适当选择 $Q_1,\mathcal L$ 时，$\zeta_-(Q_1,\mathcal L,\mathfrak s)$ 是 $\operatorname{GL}_2$ 理论中的 Euler 乘积的有限和，见 [7]。

另一方面，若 $k$ 奇，则 $D(\mathfrak s,f)$ 与 $\zeta_-(Q_1,\mathcal L,\mathfrak s)$ 不具有通常的 $\operatorname{GL}_2$ 型 Euler 乘积。不过，[18] 为半整权 $\widetilde{\operatorname{SL}}_2$ 自守形式建立了一种修正的 Euler 乘积理论。特别地，若 $f\in S_{2s}(\Gamma_0(t),\beta)$ 是 [18] 意义下的 Hecke 特征函数，则部分 Dirichlet 级数
\[
 \sum_{n\ge1}a_f(|d|n^2)n^{-\mathfrak s},
\]
其中 $d$ 为 $\Q$ 的一个虚二次扩张的判别式，具有 Euler 乘积，且在几乎所有 $p$ 处的局部因子分子为一次、分母为二次。对适当的 $\mathcal L$，也可为
$\sum_{n\ge1}M(Q_1,\mathcal L,dn^2)n^{-\mathfrak s}$ 找到类似的 Euler 乘积。于是，通过纯代数方法可证明 Rankin 卷积
\[
 \sum_{n\ge1}a_f(|d|n^2)M(Q_1,\mathcal L,dn^2)n^{-\mathfrak s}
\]
具有 Euler 乘积，且在几乎所有 $p$ 处的局部因子分子为三次、分母为四次；$k=5$ 的情形见 [16]，其中 $O(3,2)$ 与 $\operatorname{Sp}_2(\R)$ 局部同构。

\paragraph*{说明（Notes）.}
关于“双曲面分析”的完整参考文献，请参看 [21]。关于 Weil 表示与自守形式的阿代尔处理，见 [4]。

我们采用如下术语：若函数 $\phi$ 满足
\[
 (\phi|\gamma)(z)=(cz+d)^{-s}\phi(\gamma(z))
\]
具有形如
\[
 \sum_{n\ge0}c_ne^{2\pi i[(n+\kappa)/N]z}
\]
的展开，而且 $\kappa=0$ 时有 $c_0=0$，则称 $\phi$ 在尖点 $\gamma(\infty)=\mathfrak a$ 处消失。这里 $\kappa$ 是乘子在 $\mathfrak a$ 处的分歧，$N$ 是满足
\[
 \gamma\begin{pmatrix}1&N\\0&1\end{pmatrix}\gamma^{-1}\in\Gamma_1
\]
的最小正整数，$\Gamma_1$ 为所讨论的算术群，见 [18]。

\begin{thebibliography}{23}
\bibitem[1]{c19-1} T. Asai, \emph{On the Doi-Naganuma lifting associated with imaginary quadratic fields}, 预印本。
\bibitem[2]{c19-2} A. Borel, \emph{Introduction aux groupes arithmétiques}, Hermann, Paris, 1969.
\bibitem[3]{c19-3} L. Ehrenpreis, \emph{The use of partial differential equations for the study of group representations}, Proc. Sympos. Pure Math., vol. 26, Amer. Math. Soc., Providence, R.I., 1974.
\bibitem[4]{c19-4} S. Gelbart, \emph{Weil's representation and the spectrum of the metaplectic group}, Lecture Notes in Math., vol. 530, Springer-Verlag, Berlin and New York, 1976.
\bibitem[5]{c19-5} Harish-Chandra, \emph{Automorphic forms on semisimple Lie groups}, Lecture Notes in Math., vol. 62, Springer-Verlag, Berlin and New York, 1965.
\bibitem[6]{c19-6} S. Kudla, \emph{Theta-functions and Hilbert modular forms}, 预印本。
\bibitem[7]{c19-7} H. Maass, \emph{Automorphe Funktionen und indefinite quadratische Formen}, Sitzungsberichte der Heidelberger Akademie der Wissenschaften, 1949.
\bibitem[8]{c19-8} H. Maass, \emph{Verteilung der Punkte in Gittern mit indefiniter Metrik}, Math. Ann. \textbf{138} (1959), 286--315.
\bibitem[9]{c19-9} H. Maass, \emph{Über eine neue Art von nicht analytischen automorphen Funktionen und die Bestimmung Dirichletscher Reihen durch Funktionalgleichungen}, Math. Ann. \textbf{121} (1949), 141--183.
\bibitem[10]{c19-10} S. Niwa, \emph{Modular forms of half integral weight and the integral of certain theta functions}, Nagoya Math. J. \textbf{56} (1975), 147--163.
\bibitem[11]{c19-11} T. Oda, \emph{On modular forms associated with indefinite quadratic forms of signature $(2,n-2)$}, 预印本。
\bibitem[12]{c19-12} S. Rallis and G. Schiffman, \emph{Weil representation. I. Intertwining distributions and discrete spectrum}, 预印本。
\bibitem[13]{c19-13} S. Rallis and G. Schiffman, \emph{Automorphic forms constructed from the Weil representation: holomorphic case}, Amer. J. Math., 即将发表。
\bibitem[14]{c19-14} S. Rallis and G. Schiffman, \emph{Discrete spectrum of the Weil representation}, Bull. Amer. Math. Soc. \textbf{83} (1977), 267--270.
\bibitem[15]{c19-15} S. Rallis and G. Schiffman, \emph{Automorphic cusp forms constructed from the Weil representation}, Bull. Amer. Math. Soc. \textbf{83} (1977), 271--275.
\bibitem[16]{c19-16} S. Rallis and G. Schiffman, \emph{On a relation between $\widetilde{\operatorname{SL}}_2$ cusp forms and cusp forms on tube domains associated to orthogonal groups}, 预印本。
\bibitem[17]{c19-17} G. Shimura, \emph{Introduction to the arithmetic theory of automorphic functions}, Iwanami Shoten, 1971.
\bibitem[18]{c19-18} G. Shimura, \emph{On modular forms of half integral weight}, Ann. of Math. (2) \textbf{97} (1973), 440--481.
\bibitem[19]{c19-19} T. Shintani, \emph{On construction of holomorphic cusp forms of half integral weight}, Nagoya Math. J. \textbf{58} (1975), 83--126.
\bibitem[20]{c19-20} C. L. Siegel, \emph{Lectures on quadratic forms}, Tata Institute Lecture Notes, no. 7, 1955.
\bibitem[21]{c19-21} R. Strichartz, \emph{Harmonic analysis on hyperboloids}, J. Functional Analysis \textbf{12} (1973), 341--383.
\bibitem[22]{c19-22} A. Weil, \emph{Sur certaines groupes d'opérateurs unitaires}, Acta Math. \textbf{111} (1964), 143--211.
\bibitem[23]{c19-23} D. Zagier, \emph{Modular forms associated to real quadratic fields}, Invent. Math. \textbf{30} (1975), 1--48.
\end{thebibliography}
\noindent Princeton 大学（Princeton University）\\
现址：Ohio State 大学（Ohio State University）。
